\chapter{Local Workflows and Explicit Dependencies}
\label{ch:lessons-01-02}

This chapter, and the two that follow it, present Lessons~1 through~7 of the
project's tutorial sequence in full, with the discursive expansion proper to
a textbook. Each lesson retains its original pedagogical scaffold ---
objective, concept, code, expected behavior, verification, scientific
practice note --- inside a numbered \emph{Lesson box}, and is followed by
additional theoretical commentary connecting it to the architecture developed
in Part~I.

\begin{lessonbox}{Lesson 1 --- Create a Local Workflow}
\textbf{Objective.} Create and run the smallest useful DAGonStar workflow.

\textbf{Concept.} A \texttt{Workflow} owns tasks. A \texttt{DagonTask}
factory creates task instances. A local batch task executes a shell command
on the same machine.

\textbf{Code} (\texttt{lesson\_01.py}):
\begin{lstlisting}
from dagon import Workflow
from dagon.task import DagonTask, TaskType


if __name__ == "__main__":
    workflow = Workflow("Lesson01")
    task = DagonTask(TaskType.BATCH, "hello", "echo 'hello from DAGonStar' > hello.txt")

    workflow.add_task(task)
    workflow.make_dependencies()
    workflow.run()
\end{lstlisting}

\textbf{Run:}
\begin{lstlisting}[style=bash]
cp dagon.ini.sample dagon.ini
python lesson_01.py
\end{lstlisting}

\textbf{Expected behavior.} DAGonStar creates a scratch directory under the
configured \texttt{scratch\_dir\_base}, executes the task, and writes
\texttt{hello.txt} in that task's directory.

\textbf{Verify:}
\begin{lstlisting}[style=bash]
find /tmp -maxdepth 1 -type d -name '*-hello' | tail -n 1
ls /tmp/*-hello
\end{lstlisting}

\textbf{Scientific practice note.} Even a one-task workflow is useful when
the task has scientific meaning, such as downloading observations or
generating a reproducible input dataset.
\end{lessonbox}

\section{Discussion: The Smallest Useful Workflow}

Lesson~1 is deliberately trivial, but it already exercises the complete
execution lifecycle of \Cref{sec:execution-lifecycle}: a \texttt{Workflow} is
constructed, a single \texttt{BATCH} task is added, \texttt{make\_dependencies()}
parses its (dependency-free) command, \texttt{Validate\_WF()} trivially
accepts the single-node graph, and \texttt{run()} starts, executes, and
completes the one task thread. The pedagogical purpose of beginning here ---
rather than with a multi-task DAG --- is to let the scratch-directory and
\texttt{.dagon/} mechanism (formalized in \Cref{ch:lessons-03-04}) be observed
in isolation, before dependency resolution is introduced as a second source
of complexity.

\begin{lessonbox}{Lesson 2 --- Add Explicit Task Dependencies}
\textbf{Objective.} Create a workflow where tasks execute in a controlled
order.

\textbf{Concept.} Use \texttt{consumer.add\_dependency\_to(producer)} when one
task must wait for another.

\textbf{Code} (\texttt{lesson\_02.py}):
\begin{lstlisting}
from dagon import Workflow
from dagon.task import DagonTask, TaskType


if __name__ == "__main__":
    workflow = Workflow("Lesson02")

    prepare = DagonTask(TaskType.BATCH, "prepare", "echo prepared > prepared.txt")
    simulate = DagonTask(TaskType.BATCH, "simulate", "echo simulated > simulated.txt")
    summarize = DagonTask(TaskType.BATCH, "summarize", "echo summarized > summarized.txt")

    workflow.add_task(prepare)
    workflow.add_task(simulate)
    workflow.add_task(summarize)

    simulate.add_dependency_to(prepare)
    summarize.add_dependency_to(simulate)

    workflow.Validate_WF()
    workflow.run()
\end{lstlisting}

\textbf{Run:}
\begin{lstlisting}[style=bash]
python lesson_02.py
\end{lstlisting}

\textbf{Expected behavior.} The three tasks run in the order:
\texttt{prepare} $\to$ \texttt{simulate} $\to$ \texttt{summarize}.

\textbf{Verify:}
\begin{lstlisting}[style=bash]
find /tmp -maxdepth 1 -type d \( -name '*-prepare' -o -name '*-simulate' -o -name '*-summarize' \)
\end{lstlisting}

\textbf{Scientific practice note.} Explicit dependencies are best when the
relation is conceptual or procedural, for example: ``validate parameters
before launching the model.''
\end{lessonbox}

\begin{figure}[h]
\centering
\begin{tikzpicture}[
  node distance=18mm,
  task/.style={rectangle, rounded corners, draw=dagonblue, very thick,
    fill=blue!4, minimum width=30mm, minimum height=10mm, font=\sffamily\small},
  arr/.style={-{Stealth[length=2.6mm]}, thick, dagonblue}
]
  \node[task] (P) {prepare};
  \node[task, right=of P] (S) {simulate};
  \node[task, right=of S] (U) {summarize};
  \draw[arr] (P) -- node[above,font=\sffamily\scriptsize]{\texttt{add\_dependency\_to}} (S);
  \draw[arr] (S) -- node[above,font=\sffamily\scriptsize]{\texttt{add\_dependency\_to}} (U);
\end{tikzpicture}
\caption{The strictly linear, explicitly-declared dependency chain of
Lesson~2. No file is referenced between tasks; the ordering encodes a
procedural, not a data, relationship.}
\label{fig:lesson02-chain}
\end{figure}

\section{Discussion: Explicit versus Inferred Ordering}

Lesson~2 isolates the \emph{explicit} half of the task-oriented/data-oriented
distinction introduced in \Cref{sec:task-vs-data}. Note carefully what this
example does \emph{not} do: none of the three tasks reads a file produced by
another. The ordering \texttt{prepare} $\to$ \texttt{simulate} $\to$
\texttt{summarize} exists purely because the workflow author asserted it with
\texttt{add\_dependency\_to()}. This is the correct modeling choice whenever
the real-world constraint is procedural rather than a literal data
dependency --- the canonical example given in the documentation is
\emph{``validate parameters before launching the model''}: there may be no
file the validation step hands to the model step, yet the ordering is still
scientifically mandatory.

The call to \texttt{workflow.Validate\_WF()} immediately after building the
dependency graph, but before \texttt{run()}, is good practice independent of
whether dependencies are explicit or inferred: it surfaces a cycle
(\Cref{sec:dag-theory}) as an explicit, early modeling error rather than as a
runtime deadlock.

\section{Mixing Explicit and Inferred Dependencies in One Workflow}

Lessons~1 and~2 are presented separately for pedagogical clarity, but it is
worth being explicit, before moving on, that nothing in the engine restricts
a single workflow to using only one dependency-expression style. A realistic
pipeline routinely mixes both: a \texttt{validate\_parameters} task that
gates \texttt{run\_model} through an explicit \texttt{add\_dependency\_to()}
call (because the relationship is procedural, as Lesson~2 demonstrates), and
a \texttt{run\_model} task that simultaneously consumes
\texttt{prepare\_input}'s output file through a \texttt{workflow://}
reference (because that relationship is a genuine data dependency, as
Lesson~3 will demonstrate in the next chapter).

\begin{example}
Consider a workflow with four tasks: \texttt{check\_license} (verifies that a
required model license file is present and current --- a procedural
precondition with no file product consumed by \texttt{run\_model});
\texttt{prepare\_input} (writes \texttt{forcing.nc}); \texttt{run\_model}
(consumes \texttt{forcing.nc} via \texttt{workflow://} \emph{and} must wait
for \texttt{check\_license}, expressed via \texttt{add\_dependency\_to()});
and \texttt{archive} (consumes \texttt{run\_model}'s output via
\texttt{workflow://}). The resulting dependency graph has three edges of two
different kinds: one explicit ($\texttt{check\_license} \to
\texttt{run\_model}$) and two inferred ($\texttt{prepare\_input} \to
\texttt{run\_model}$, $\texttt{run\_model} \to \texttt{archive}$). Both
kinds are recorded identically in \texttt{prevs}/\texttt{nexts}
(\Cref{ch:architecture}) once \texttt{make\_dependencies()} has run, and
\texttt{Validate\_WF()} treats the resulting four-node graph as a single,
uniform DAG regardless of how each edge was created.
\end{example}

\begin{practicebox}
A useful authoring discipline, consistent with the modeling guidance of
\Cref{sec:task-vs-data}: when writing a new task, first ask whether the task
\emph{reads a file} that another task produces. If so, express that
dependency through \texttt{workflow://} and let it be inferred --- doing so
automatically keeps the declared dependency synchronized with the actual
data the command consumes, as \Cref{ch:schema} discusses at length. Reserve
explicit \texttt{add\_dependency\_to()} calls for the residual cases where no
file changes hands, but an ordering constraint is nonetheless real.
\end{practicebox}

\section{Validating the Graph Without Running It}

Lesson~2's use of \texttt{Validate\_WF()} before \texttt{run()} previews a
pattern this book will return to repeatedly in Part~II, under the name
\emph{local structural verification} (\Cref{ch:backends}): separating the
act of \emph{constructing and validating} a dependency graph from the act of
\emph{executing} it. This separation has a concrete pedagogical and
debugging benefit even at this early stage of the tutorial sequence. A
workflow author uncertain whether a hand-edited set of
\texttt{add\_dependency\_to()} calls has accidentally introduced a cycle can
call \texttt{make\_dependencies()} followed by \texttt{Validate\_WF()} alone,
inspect the raised exception (or its absence) and the resulting
\texttt{prevs}/\texttt{nexts} lists on each task object, and only then decide
whether to proceed to \texttt{run()} --- without incurring the cost, or the
side effects, of actually executing any task's command. For workflows whose
tasks have expensive or non-idempotent side effects (a remote job
submission, a billable cloud operation, an API call with rate limits), this
validate-before-run discipline is not merely convenient but essential.

\begin{exercise}
Rewrite Lesson~2 so that \texttt{simulate} actually consumes the file written
by \texttt{prepare} (rename \texttt{prepared.txt} to a path inside an
\texttt{output/} subdirectory), and replace the explicit
\texttt{add\_dependency\_to()} calls with \texttt{workflow://} references.
Confirm that \texttt{workflow.make\_dependencies()} reconstructs the same
linear order automatically. This exercise previews \Cref{ch:lessons-03-04}.
\end{exercise}

\begin{exercise}
Extend the four-task example of this section's worked example into a
runnable DAGonStar script. Verify, using
\texttt{task.prevs}/\texttt{task.nexts} inspection as suggested above, that
the engine records exactly three dependency edges, and confirm that
\texttt{Validate\_WF()} accepts the graph without raising an exception.
\end{exercise}
